\begin{frame}
  \frametitle{Automation}
  \begin{figure}[htpb]
      \centering
      \includegraphics[width=2cm]{images/robot-icon.eps}
  \end{figure}
  \begin{center}
  {\tiny Source: \cite{robot_icon}}
  \end{center}

  Automating out repetive tasks saves more time for designing and developing features and bug fixes.
  \newline
  \newline
  Question: How do we implement automation?
  \pause
  \newline
  \newline
  Answer: Using an {\it automation framework}.

\end{frame}

\begin{frame}[t]
    \frametitle{Automation Framework}
    \begin{figure}[htpb]
        \begin{subfigure}
            \centering
            \includegraphics[width=3cm]{images/circleci-logo.eps}
        \end{subfigure}
        \begin{subfigure}
            \centering
            \includegraphics[width=3cm]{images/github-actions-logo.png}
        \end{subfigure}
    \end{figure}
    \begin{center}
        {\tiny Sources: \cite{circle_ci_logo}, \cite{github_actions_logo}}
    \end{center}

    Automation frameworks are services execute user-created instruction sets called {\it workflows} when certain conditions, or {\it triggers}, are met.

    Automation frameworks can read files from the main host repository. This enables users to create workflows to do things like:
    \begin{itemize}
        \item automatically run a test suite whenever the code changes ({\it continuous integration})
        \item build and deploy online documentation
        \item manage repository metadata
    \end{itemize}
    
\end{frame}

\begin{frame}
    \frametitle{GitHub Actions}
    
    GitHub Actions is an automation framework integrated into every GitHub repository.

    \begin{figure}[htpb]
        \centering
        \includegraphics[width=10cm]{images/overview-actions-simple}
        \newline
        {\tiny Source: \cite{github_actions_conceptual_example}}
    \end{figure}

    The basic workflow file structure is as follows:
    \begin{itemize}
        \item Workflow name
        \item Define workflow triggering events
        \item Define the workflow jobs and steps
    \end{itemize}

\end{frame}

\begin{frame}[fragile]
    \frametitle{GitHub Actions Workflow Example}
    \framesubtitle{Populating SaltProc release notes}

    \begin{minted}{yaml}
# Preamble
name: Populate SaltProc release notes

on:
  push:
    branches:
      - 'master'
    paths:
      - 'doc/releasenotes/v**.rst'
  # enable worflow to be run manually
  workflow_dispatch:

jobs:
  populate-releasenotes:
    runs-on: ubuntu-latest
    defaults:
      run:
        shell: bash -l {0}
    \end{minted}
    
    

    %\begin{figure}[htpb]
    %    \centering
    %    \includegraphics[width=10cm]{images/github_actions_workflow_ex1}
    %\end{figure}

\end{frame}

\begin{frame}[fragile]
    \frametitle{GitHub Actions Workflow Example}
    \framesubtitle{Populating SaltProc release notes}
    \begin{minted}{yaml}
    steps:
      - uses: actions/checkout@v2

      - name: Set up Python 3.9
        uses: actions/setup-python@v2
        with:
          python-version: 3.9
      - name: Add conda to system path
        run: |
          # $CONDA is an environment variable pointing to the
          # root of the miniconda directory
          echo $CONDA/bin >> $GITHUB_PATH

      - name: install pandoc 
        run: |
          conda install -c conda-forge pandoc
          pip install --upgrade pandoc
    \end{minted}

    %\begin{figure}[htpb]
    %    \centering
    %    \includegraphics[width=10cm]{images/github_actions_workflow_ex2}
    %\end{figure}
\end{frame}

\begin{frame}[fragile]
    \frametitle{GitHub Actions Workflow Example}
    \framesubtitle{Populating SaltProc release notes}
    \begin{minted}{yaml}
      - name: Get most recent draft release version
        run: |
          echo "RELEASE_VERSION=$(gh api repos/
          ${{ github.repository }}/
          releases --jq '.[0] | .name')" >> $GITHUB_ENV
          echo "RELEASE_ID=$(gh api repos/
          ${{ github.repository }}/
          releases --jq '.[0] | .id')" >> $GITHUB_ENV
        env:
          GITHUB_TOKEN: ${{ secrets.GITHUB_TOKEN }}

      - name: Convert .rst to .md 
        run: |
          pandoc -o RELEASENOTES.md -f rst -t gfm doc/releasenotes/
          ${{ env.RELEASE_VERSION }}.rst --columns 1000
          sed -i "s/# Release notes for ${{ env.RELEASE_VERSION }}//g"
          RELEASENOTES.md
    \end{minted}

    %\begin{figure}[htpb]
    %    \centering
    %    \includegraphics[width=10cm]{images/github_actions_workflow_ex2}
    %\end{figure}
\end{frame}


\begin{frame}[fragile]
    \frametitle{GitHub Actions Workflow Example}
    \framesubtitle{Populating SaltProc release notes}

    \begin{minted}{yaml}
      - name: Populate the release description with RELEASENOTES.md
        run: |
          CURRENT_TAG=$(gh api repos/${{ github.repository }}/
          releases/${{ env.RELEASE_ID }} \
          -H "Authorize: token ${{ secrets.GITHUB_TOKEN }}" \
          -H "Accept: application/vnd.github.v3+json" \
          -X GET \
          --jq '.tag_name')
          gh api repos/${{ github.repository }}/releases/
          ${{ env.RELEASE_ID }} \
          -H "Authorize: token ${{ secrets.GITHUB_TOKEN }}" \
          -H "Accept: application/vnd.github.v3+json" \
          -X PATCH \
          -F tag_name=$CURRENT_TAG \
          -F body="$(cat RELEASENOTES.md)"
        env:
          GITHUB_TOKEN: ${{ secrets.GITHUB_TOKEN }}        
    \end{minted}
    %\begin{figure}[htpb]
    %    \centering
    %    \includegraphics[width=10cm]{images/github_actions_workflow_ex3}
    %\end{figure}
\end{frame}


